\documentclass[11pt]{article}
\usepackage[a4paper,margin=1in]{geometry}
\usepackage{amsmath,amssymb}
\usepackage{booktabs}
\usepackage{array}
\usepackage[numbers,sort&compress]{natbib}
\usepackage{hyperref}

\newcommand{\derived}{\textit{[derived here]}}
\newcommand{\fromSM}{\textit{[standard SM]}}
\newcommand{\fromISPG}[1]{\textit{[from #1]}}
\newcommand{\numerical}{\textit{[numerically verified]}}
\newcommand{\deferred}{\textit{[deferred to extension program]}}
\newcommand{\modeldep}{\textit{[model-dependent]}}

\title{Structural Compatibility of the Standard Model with\\
Induced Scalar-Pressure Gravity:\\
a Partial Embedding and Extension Program}
\author{George Lotishvili}
\date{Version 1.0 (Companion appendix to the MASS paper)}

\begin{document}
\maketitle

\begin{abstract}
We present a six-phase analysis of how the Standard Model (SM)
gauge and fermion content can be hosted in the cavity framework
of Induced Scalar-Pressure Gravity (ISPG). The scope is
\emph{compatibility}, not derivation from first principles: given
the ISPG archion substrate and the Mathieu ladder previously
established, we check whether the SM quantum numbers, anomaly
conditions, and electroweak structure can be assigned
consistently to cavity labels, and which elements of the SM must
still be imported. We find:
(i) the count of confined modes of the $D_3$-deformed oscillon
is eight, matching $\dim\mathrm{SU}(3)_{\mathrm{adj}}$, without
fine-tuning; the $\mathrm{SU}(3)$ Lie-algebra structure itself
is not independently derived from the cavity dynamics in this
work;
(ii) the integer cavity labels $N_W,N_Z$ individually fit
$m_W,m_Z$ to $\leq 0.31\%$; their ratio $(N_W/N_Z)^2=0.8817$
matches the on-shell mass-ratio convention
$\cos\theta_W^{\rm os}\equiv m_W/m_Z\simeq0.8815$ to
$0.03\%$---an
\emph{internal consistency check}, not an independent
prediction of the Weinberg angle, since the two indices are
themselves fits to the masses whose ratio defines
$\cos\theta_W^{\rm os}$ at SM tree level;
(iii) the ISPG action contains no $\Phi^4$ term, so a SM-type
Higgs mechanism is ruled out at the fundamental level; the
specific replacement dynamics for gauge-boson mass generation
is not derived here and is listed as an extension item;
(iv) the hypercharge pattern $Y\cdot N_c\in\mathbb{Z}$ is a
compatibility and integer-grading clue respected by the imported
SM assignment; it is not independently derived here, and the
SM-specific hypercharge values are imported via the
generation-1 fitting formula plus Gell-Mann--Nishijima;
(v) the three generations are interpreted (under a
model-dependent three-axis hypothesis) as the three orthogonal
vibration axes of the 3D archion, giving a geometric reading of
the number $3$;
(vi) neutrinos lie \emph{outside} the Mathieu-ladder fermion
formula and are taken as polarization-rotation modes per the
companion MASS paper; their inclusion in the formula would
mispredict $m_\nu$ by six orders of magnitude, and the
exclusion is therefore structural, not cosmetic.
Extension topics (the 3D vector cavity spectrum behind the
individual indices $N_W,N_Z$; the $\mathrm{SU}(3)$ Lie-algebra
structure beyond the count; the archion-interior mechanism for
gauge-boson mass generation; the fine-structure constant;
CKM and PMNS matrices; the exact charged-lepton mass
hierarchy) are identified as a programmatic continuation and
do not conflict with the compatibility established here.
\end{abstract}

\tableofcontents

\section{Introduction and scope}
\label{sec:intro}

The ISPG framework
\cite{Lotishvili2026,Lotishvili2026Q} postulates a single real scalar
field $\Phi(x)$ coupled minimally to gravity through a pressure term.
Earlier ISPG papers \cite{Lotishvili2026,Lotishvili2026Q} establish:
(a) the bare action contains only the curvature $R$ and the kinetic
term $\tfrac{1}{2}(\partial\Phi)^2$, with \emph{no} non-derivative
potential; (b) localized finite-energy solutions (``archions'') exist
as time-periodic oscillons inside a finite cavity; (c) the linear
spectrum of small fluctuations around the archion obeys Mathieu's
equation, whose characteristic values $b_N(q)$ at fixed
$q\simeq 1.853$ reproduce all known charged elementary particle
masses through $m_N = m_1\,b_N(q)$ to better than $0.3\%$
\cite{McLachlan1947,DLMFMathieu,PDG2024}.

\paragraph{Gap addressed.}
What the earlier ISPG papers did \emph{not} settle was whether
SM quantum numbers, anomaly conditions, and electroweak
structure can be \emph{hosted consistently} on the cavity
labels that ISPG natively provides---and, if so, which SM
ingredients are imported from external knowledge versus derived
from the cavity. This paper carries out that compatibility
check explicitly across six phases, with a clear separation
(via epistemic labels) of what is derived, what is fit, what is
imported, and what is deferred. The result is a
\emph{partial} structural embedding: sufficient to exclude gross
inconsistencies and to identify the specific open items
(§\ref{sec:extension}) that a \emph{complete} embedding would
still need to resolve.

\paragraph{Methodological commitments.}
Throughout, each physical statement carries an explicit
epistemic label:
\derived\ (proved in this paper from the ISPG substrate),
\fromSM\ (used as a known SM fact),
\numerical\ (verified numerically by a Python script in
\texttt{python/}),
\modeldep\ (holds under a specific hypothesis, stated as such), or
\deferred\ (part of the extension program). The six phases A--F
below partition the embedding into structurally separable tasks;
each phase terminates in an explicit pass/fail criterion.

\paragraph{Non-goals.}
This paper does \emph{not} derive the exact masses of the muon and
tau, does \emph{not} derive the Cabibbo--Kobayashi--Maskawa (CKM)
or Pontecorvo--Maki--Nakagawa--Sakata (PMNS) mixing matrices, and
does \emph{not} reconstruct the fine-structure constant $\alpha$
from first principles. These are open targets for the extension
program (\S\ref{sec:extension}). Their \emph{absence} does not
invalidate the structural embedding; it constrains the scope of
what the embedding establishes.

\section{ISPG prerequisites used below}
\label{sec:prereq}

The following ISPG results are taken as inputs to the embedding;
each is proved in the referenced paper or Python verification
script.

\paragraph{Action.} \fromISPG{\cite{Lotishvili2026}}
The bare action is
\begin{equation}
  \label{eq:action}
  S = \frac{1}{16\pi G}\int d^4x\,\sqrt{-g}\,
  \left[\,R + \tfrac{1}{2}(\partial\Phi)^2\,\right] + S_m,
\end{equation}
with $S_m$ the standard matter action. There is no $\Phi^2$,
$\Phi^4$, or generic $V(\Phi)$ term; consequently any potential
appearing below is \emph{effective}, induced by nonlinear
self-interaction through archion dynamics
\cite{Copeland1995,GleiserWatkins1993}, not added by hand.

\paragraph{Mathieu ladder.} \fromISPG{\cite{Lotishvili2026}}
Linearized fluctuations around the periodic archion solve
Mathieu's equation. At the characteristic deformation
$q\simeq 1.853$, the characteristic values $b_N(q)$ form a
discrete ladder indexed by $N=1,2,3,\ldots$, and every massive SM
particle sits on an integer $N$. The mass law
\begin{equation}
  \label{eq:massN}
  m_N = m_1 \cdot b_N(q)
\end{equation}
is the master interpretation of $N$ as a cavity-mode label.
Empirical fits (for the electron taken as $N_e=5$) give
$m_N\propto N^2$ to high accuracy in the heavy sector; this
quadratic scaling is used in phases D and F
(\texttt{phase\_d\_ew\_breaking.py}).

\paragraph{Archion and $D_3$.} \fromISPG{\cite{Lotishvili2026Q}}
An archion with triangular ($D_3$) deformation in three spatial
directions admits a tensor-mode spectrum whose confined sector
contains exactly eight states for a robust range of the
deformation parameter (documented in
\texttt{SM/ISPG\_EightGluon\_Robustness.md});
this is used in \S\ref{sec:phaseA}.

\paragraph{Three-axis hypothesis.} \numerical\ \modeldep\
The 3D archion can sustain three orthogonal vibration modes
(one per spatial axis), each with its own damping. Under this
hypothesis, the $(e,\mu,\tau)$ charged leptons are identified
with one oscillator on three axes, and the Koide relation is
reproduced to $0.0006\%$, Sargent's law to $0.26\%$, the
absolute $\tau$ mass to $0.4\%$
(session of 2026-04-21, see \texttt{MASS/WORKING\_N\_origin\_and\_59.md}
\S\S\,1--10). This hypothesis is model-dependent; the
embedding in \S\ref{sec:phaseF} uses it to \emph{interpret}
the number $3$ of generations but does not need the exact
numerical reproduction of the three charged-lepton masses.

\paragraph{Two forms of the mass law.}
The $0.0006\%$ Koide figure quoted above is obtained with the
full Mathieu characteristic values $a_N(q\!\simeq\!1.853)$ from
the linearized archion spectrum. The large-$N$ approximation
$m\propto N^2$, used elsewhere in this paper for the heavy
sector (electroweak bosons, Higgs, and the cross-sector
compatibility survey), is the asymptotic limit of the same
ladder: it matches the observed masses to sub-percent precision
for $N\gtrsim 70$ but returns $Q=\tfrac12$ for Koide rather than
the observed $Q\simeq\tfrac23$. Precision-level lepton
statements (Koide, Sargent) therefore use the full Mathieu
evaluation, while compatibility-level embedding statements in
\S\S\ref{sec:phaseA}--\ref{sec:phaseF} use the $N^2$ form. No
additional physical input is involved; the two forms agree
asymptotically.

\section{Phase A --- Gauge sector}
\label{sec:phaseA}

\paragraph{Script.} \texttt{python/phase\_a\_gauge\_sector.py}.

\subsection{$\mathrm{SU}(3)$: eight confined gluons from $D_3$}
\label{sec:phaseA_su3}

\numerical\ (octet-count clue; documented in the robustness
analysis, supplementary notes on file with the author; summarized
here).
The $D_3$-deformed three-dimensional oscillon carries a tensor
perturbation spectrum with ten modes at leading order; exactly
eight are confined (lie below the chosen leak threshold) across a
nonzero, threshold-sensitive region of parameter space
\begin{equation}
  g_0 \in [0.1, 8.0], \qquad V \in [3.0, 9.5],
\end{equation}
with the central value $V=\Delta\approx 6.26$ identified as the
resonant deformation. The confined multiplicity is
\emph{discrete-binary} in $\{8,10\}$: the values $\{5,6,7,9\}$
never appear. This matches the dimension of the
$\mathrm{SU}(3)$ adjoint representation ($\dim=8$) as a structural
compatibility clue. The robustness analysis gives an $N=8$ region
of $13.23\%$ at the default $40\%$ leak threshold, $6.50\%$ at
$45\%$, and $0\%$ at $50\%$; the physical threshold and the
parameters $g_0,V$ must still be derived from the full
$G_2(X)$ dynamics.

\paragraph{Caveat: count versus algebra.}
\deferred\ The count of eight confined modes is a necessary but
not sufficient condition for an $\mathrm{SU}(3)$ gauge
structure. The full identification requires that the eight
modes satisfy the Gell-Mann commutation relations
$[T^a,T^b]=if^{abc}T^c$ with the standard $\mathrm{SU}(3)$
structure constants $f^{abc}$. Extracting these constants from
the nonlinear archion dynamics (specifically, from the
trilinear mode-coupling vertices) is a separate computational
task deferred to item E6 of the extension program. In its
current form, the result of this subsection is therefore
``$\dim\mathrm{SU}(3)_{\mathrm{adj}}$ is natively reproduced,''
not ``$\mathrm{SU}(3)$ is derived.''

\subsection{$\mathrm{SU}(2)\times\mathrm{U}(1)_Y$: numerical
consistency}
\label{sec:phaseA_ew}

\numerical\ The electroweak bosons sit on specific Mathieu
indices:
\begin{center}
\begin{tabular}{lcccc}
\toprule
Particle & $m$ (GeV) & $N$ & $m_{\text{pred}}=(N/5)^2 m_e$ &
  error \\
\midrule
$W^\pm$ & $80.369$ & $1986$ & $80.62$\,GeV & $0.31\%$ \\
$Z^0$  & $91.188$ & $2115$ & $91.44$\,GeV & $0.27\%$ \\
$H^0$  & $125.250$& $2479$ & $125.61$\,GeV& $0.29\%$ \\
\bottomrule
\end{tabular}
\end{center}
\noindent
This table uses the rough large-$N$ estimate
$m_N=(N/5)^2m_e$.  The full Mathieu normalization gives a
separate precision diagnostic
($W\simeq80.389$~GeV, $Z\simeq91.172$~GeV,
$H\simeq125.254$~GeV), but it still does not derive the
labels $(N_W,N_Z,N_H)$.
\numerical\ The ratio
\begin{equation}
  \label{eq:cosThetaW}
  \left(\frac{N_W}{N_Z}\right)^{\!2}
  \;=\;\left(\frac{1986}{2115}\right)^{\!2}
  \;=\;0.8817,
\end{equation}
is to be compared with the on-shell mass-ratio convention
$\cos\theta_W^{\rm os}\equiv m_W/m_Z\simeq0.8815$
\cite{PDG2024} and with the observed
$m_W/m_Z=0.8814$: a $\leq 0.03\%$ agreement in each case.

\paragraph{Nature of this test: consistency, not independent
prediction.}
Since $N_W$ and $N_Z$ were themselves fit to $m_W$ and $m_Z$
individually (to $0.31\%$ and $0.27\%$ respectively), their
ratio \emph{automatically} matches $m_W/m_Z$ within the
combined fit residuals. The ratio test therefore does not
constitute an independent prediction of the Weinberg angle; at SM
tree level in the on-shell convention
$m_W/m_Z=\cos\theta_W^{\rm os}$ holds identically and any
theory that reproduces the two masses to high accuracy will
reproduce this ratio. The non-trivial content of
\eqref{eq:cosThetaW} is that the \emph{discrete integer}
Mathieu ladder (step size $\Delta N=1$ out of $N\sim 2000$)
can simultaneously host both bosons on the same mass law
$m\propto N^2$ with residuals well below one ladder step; the
ratio inherits this precision mechanically.

\paragraph{What is open.} \deferred\ The first-principles
derivation of the individual labels $(N_W,N_Z)$ from the
three-dimensional vector cavity spectrum is the principal
outstanding task for a non-trivial electroweak prediction
(item E1). Until it is completed, the electroweak sector of
the present embedding is a consistency check of the Mathieu
ansatz against two data points, not a structural derivation
of the Weinberg angle.

\subsection{Phase A verdict}

Yellow on the gluon count: the $D_3$-deformed oscillon gives an
octet-count compatibility clue, with partial numerical robustness
and threshold sensitivity, while the full Lie-algebra structure is
deferred (E6). Yellow on
$\mathrm{SU}(2)\times\mathrm{U}(1)_Y$: the Mathieu ladder
accommodates $W,Z,H$ with sub-percent residuals, but the
electroweak test is internal consistency rather than an
independent prediction of the Weinberg angle. The derivation of the
specific labels $(N_W,N_Z)$ is deferred (E1). No contradiction
with ISPG's single-scalar action is introduced by either step.

\section{Phase B/C --- Fermion quantum numbers}
\label{sec:phaseBC}

\paragraph{Scripts.} \texttt{python/phase\_c\_quantum\_number\_map.py},
\texttt{python/phase\_c\_step2\_T3\_Y\_map.py}.

\subsection{Generation-1 SM-to-cavity compatibility map}
\label{sec:phaseC_gen1}

\numerical\ We test whether the imported SM labels $(Q,N_c)$
can be used to assign cavity-compatible generation-1 Mathieu
indices for the massive charged fermions $(e,u,d)$. An ansatz
using these imported labels
\begin{equation}
  \label{eq:gen1formula}
  N \;=\; 5\bigl[\,1 + N_c\bigl(1 - |Q|\bigr)\bigr]
\end{equation}
reproduces the three observed indices. The implied structure is
\begin{equation}
  3|Q| \;\in\;\{0,1,2,3\},
\end{equation}
i.e. $Q$ takes only values whose triple is an integer: this is a
third-integer charge grading. Interpreting this grading as a
$\mathbb{Z}_3$ or $D_3$ representation requires a separate
representation map and is left open.

\paragraph{Status: fit, not prediction.}
The formula~\eqref{eq:gen1formula} has three free ingredients
(the anchor $5=N_e$, the coefficient $1$ multiplying $N_c$, and
the choice of functional form $1-|Q|$) and is fit to three data
points $(N_e,N_u,N_d)$. A perfect match is therefore guaranteed
by parameter counting and is not by itself predictive. The
structural content that \emph{does} survive parameter counting
is the \emph{shape}: that a formula built from the imported SM
labels $(Q,N_c)$, anchored on the electron and returning
cavity-compatible integer multiples of $5$, exists at all. A
genuine prediction from~\eqref{eq:gen1formula} would require
either (i) a fourth data point that the formula was not fit to,
or (ii) an independent derivation of the coefficients from the
archion dynamics. Both are deferred.

\subsection{Neutrinos are excluded from the Gen-1 formula}
\label{sec:neutrinos_exclusion}

\derived\ Formally, substituting $Q=0,\,N_c=1$
into~\eqref{eq:gen1formula} gives $N_{\nu_e}=10$, which
via the $m\propto N^2$ scaling would imply
$m_{\nu_e}\approx (10/5)^2 m_e=2.044\,\text{MeV}$.
This exceeds the observed neutrino mass bound
$m_\nu<1\,\text{eV}$ by more than six orders of magnitude.
The formula therefore \emph{cannot} apply to neutrinos.

The resolution is structural, not ad hoc: the companion MASS
paper \cite{Lotishvili2026} identifies neutrinos as
\emph{polarization-rotation} (helical-torsional) modes of
$\Phi$, distinct from the Mathieu-ladder oscillon family, with
their tiny mass arising from the suppressed torsional rigidity
of the nonlinear $G_2(X)$ coupling. Under that identification,
neutrinos are not cavity oscillons at all and the Mathieu
mass law $m\propto N^2$ does not apply to them. Consequently
the scope of~\eqref{eq:gen1formula} is explicitly restricted
to the \emph{massive charged} fermion sector, and the nominal
$N=10$ assignment for $\nu_e$ is discarded.

This exclusion eliminates a potential internal contradiction
of the embedding but, equally, limits its reach: the Gen-1
formula says nothing about neutrino quantum numbers, and the
$T_3,Y$ assignments for $\nu_L$ in Table~\ref{tab:gen1} are
imported from the SM via Gell-Mann--Nishijima, not derived
from the cavity.

\subsection{Hypercharge map and Gell-Mann--Nishijima}
\label{sec:phaseC_T3Y}

\derived\ The full left-Weyl content of generation 1 is listed in
Table~\ref{tab:gen1}. For every entry, the Gell-Mann--Nishijima
relation
\begin{equation}
  \label{eq:GMN}
  Q \;=\; T_3 + \tfrac{1}{2}Y
\end{equation}
is satisfied as an identity.

\begin{table}[h]
\centering
\caption{Generation-1 Weyl-fermion content (left-handed basis,
right-handed fields written as charge-conjugated left-handed
spinors).}
\label{tab:gen1}
\begin{tabular}{lcccccc}
\toprule
Field & $Q$ & $T_3$ & $Y$ & $N_c$ & $Y\cdot N_c$ & $Q=T_3+Y/2$\\
\midrule
$\nu_{eL}$ &   $0$  & $+\tfrac12$ & $-1$ & $1$ & $-1$ & OK\\
$e_L$      &  $-1$  & $-\tfrac12$ & $-1$ & $1$ & $-1$ & OK\\
$e_L^c$    &  $+1$  & $0$         & $+2$ & $1$ & $+2$ & OK\\
$u_L$      & $+\tfrac23$ & $+\tfrac12$ & $+\tfrac13$ & $3$ & $+1$ & OK\\
$d_L$      & $-\tfrac13$ & $-\tfrac12$ & $+\tfrac13$ & $3$ & $+1$ & OK\\
$u_L^c$    & $-\tfrac23$ & $0$         & $-\tfrac43$ & $3$ & $-4$ & OK\\
$d_L^c$    & $+\tfrac13$ & $0$         & $+\tfrac23$ & $3$ & $+2$ & OK\\
\bottomrule
\end{tabular}
\end{table}

\paragraph{Integer hypercharge-times-color.} \fromSM\ \numerical\ A direct
reading of the last-but-one column gives
\begin{equation}
  \label{eq:YNc_integer}
  Y\cdot N_c \;\in\; \{-4,\,-2,\,-1,\,+1,\,+2\}\;\subset\;\mathbb{Z}
\end{equation}
for every field (including both chiralities). This integrality
is an observed compatibility clue: the imported SM assignment
respects an integer grading when hypercharge is multiplied by
color multiplicity. It is not, by itself, a derivation of the
hypercharges or of anomaly cancellation. The anomaly equations
require the specific set $\{-4,-2,-1,+1,+2\}$ with the specific
multiplicities shown, not merely arbitrary integers. In this
work the $Y$ values are imported through the Gen-1 formula
\eqref{eq:gen1formula} combined with Gell-Mann--Nishijima; the
product $Y\,N_c$ is therefore a checked compatibility pattern,
not an automatic consequence of independently derived cavity
labels. The integrality is also a hallmark of
Pati--Salam-type unifications
(J.\,C.~Pati and A.~Salam, \emph{Phys.\ Rev.\ D} \textbf{10},
275 (1974)), reached here through a different route.

\paragraph{Chirality.} \fromSM\ \deferred\ The distinction
$(L/R)$ in the fermion content is imported from the standard
Dirac structure of $S_m$; deriving chirality from cavity
boundary conditions is a separate open problem
(\S\ref{sec:extension}).

\subsection{Phase B/C verdict}

For the massive charged generation-1 fermions $(e,u,d)$ the
indices $(N_e,N_u,N_d)$ are accommodated by a three-ingredient
fit formula~\eqref{eq:gen1formula} built from imported SM
labels $(Q,N_c)$; the fit is exact on the three data
points by parameter counting and is not yet predictive.
Neutrinos are excluded from the formula by construction
(\S\ref{sec:neutrinos_exclusion}). Gell-Mann--Nishijima holds
exactly for every entry in Table~\ref{tab:gen1} because the
$T_3,Y$ assignments are imported from the SM consistent with
it. The integer structure $Y\,N_c\in\mathbb{Z}$ is a checked
compatibility pattern respected by the imported SM assignment.

\section{Phase D --- Electroweak symmetry breaking without Higgs}
\label{sec:phaseD}

\paragraph{Script.} \texttt{python/phase\_d\_ew\_breaking.py}.

\subsection{No $\Phi^4$ in the action}

\derived\ The bare action~\eqref{eq:action} contains no
$\Phi^2$ and no $\Phi^4$ term. A SM-type Higgs mechanism is
therefore absent at the fundamental level: there is no
$V_{\mathrm{Higgs}}(\Phi)=\mu^2\Phi^2+\lambda\Phi^4$ with
$\mu^2<0$, no $\mathrm{SU}(2)$-doublet scalar, and no vacuum
expectation value $\langle\Phi\rangle=v/\sqrt{2}$ filling the
universe homogeneously. This is a structural consequence of the
single-scalar postulate and forces ISPG to generate
gauge-boson masses by a different mechanism (see below).

\subsection{Local EW breaking: the archion interior}

\modeldep\ The role that the Higgs plays in the SM---providing
a region of space where $\Phi\neq 0$ so that $W,Z$ acquire mass---
is played in ISPG, under this proposal, by the \emph{interior of
each archion}. An archion is a localized finite-energy excitation
with $\Phi(r)\neq 0$ inside its core and $\Phi(r)\to 0$ outside.
The SM's ``broken phase'' maps to the archion interior; the SM's
``unbroken phase'' maps to the laboratory vacuum. Hence EW
symmetry breaking in ISPG is proposed to be \emph{local}, not
global.

\paragraph{What is not shown.} \deferred\ The specific dynamical
mechanism by which gauge bosons acquire their observed masses
from the archion interior---the analogue of the Higgs Yukawa
couplings and gauge-kinetic term
$(D_\mu\Phi)^\dagger(D^\mu\Phi)$---is \emph{not} derived in
this paper. A complete replacement would need to show how an
$\mathrm{SU}(2)$ gauge field minimally coupled to the archion
acquires a mass matrix matching $M_W,M_Z$ (with the photon
remaining massless) from the archion's spatial profile alone.
Until this is done, the statement ``ISPG does not need a Higgs''
should be read as: ``the ISPG action forbids a SM-type
fundamental Higgs scalar by construction, and an alternative
local mechanism is proposed but not derived.'' This task is
item E7 in the extension program.

Two immediate consequences:
\begin{itemize}
  \item[(i)] The laboratory vacuum is genuinely empty of any
  Higgs condensate, consistent with the observed cosmological
  constant problem being an order-of-magnitude mismatch between
  SM-expected vacuum energy and observation.
  \item[(ii)] The apparent scale $v\approx 246$\,GeV is an
  \emph{internal length scale} of the archion, not a property of
  space at large. Via~\eqref{eq:massN} this length maps to a
  large Mathieu index ($N\approx 3471$ in the $m\propto N^2$
  extrapolation); it is not itself a particle.
\end{itemize}

\subsection{Consistency check:
$m_W/m_Z=\cos\theta_W^{\rm os}$}

\numerical\ As discussed in \S\ref{sec:phaseA_ew}, the
ratio~\eqref{eq:cosThetaW} matches the on-shell mass-ratio
convention $\cos\theta_W^{\rm os}$ to $0.03\%$. This is a
self-consistency check of the Mathieu
ansatz (the same integer ladder hosts both bosons), not an
independent prediction of the Weinberg angle, since the individual
indices $N_W,N_Z$ are themselves fits to $m_W,m_Z$. What is
genuinely non-trivial is that a \emph{discrete} ladder can
host two particles of such close but distinct masses
simultaneously with sub-percent residuals: this constrains
the ladder spacing to be fine enough at large $N$, which the
$m\propto N^2$ law satisfies automatically.

\subsection{Phase D verdict}

The ISPG action~\eqref{eq:action} forbids a SM-type fundamental
Higgs scalar by construction (no $\Phi^2$, no $\Phi^4$). A
local alternative (gauge-boson mass generation inside the
archion) is proposed but not derived (E7). The $m_W/m_Z$ ratio
consistency is a fit diagnostic, not a derivation of the
Weinberg angle. No internal contradiction with ISPG is
introduced, but the electroweak sector remains structurally
incomplete until E1 and E7 are carried out.

\section{Phase E --- Anomaly cancellation}
\label{sec:phaseE}

\paragraph{Script.} \texttt{python/phase\_e\_anomalies.py}
(exact arithmetic with \texttt{fractions.Fraction}).

\subsection{The five conditions}

\fromSM\ A consistent gauge theory with the SM matter content
must satisfy five independent anomaly-cancellation conditions:

\begin{enumerate}
  \item $\mathrm{U}(1)_Y^3$: $\sum_L Y^3 = 0$;
  \item $\mathrm{SU}(2)^2\times\mathrm{U}(1)_Y$: $\sum_{\text{doublets}} Y = 0$;
  \item $\mathrm{SU}(3)^2\times\mathrm{U}(1)_Y$: $\sum_{\text{colored}} Y = 0$
    (summed per $\mathrm{SU}(2)$ component);
  \item gravitational $\times\,\mathrm{U}(1)_Y$: $\sum_L Y = 0$;
  \item Witten $\mathrm{SU}(2)$: the total number of
    $\mathrm{SU}(2)$ doublets (counting color multiplicity) is
    even.
\end{enumerate}
Here $\sum_L$ denotes a sum over all left-handed Weyl fermions
(right-handed fields are charge-conjugated into their left-handed
equivalents).

\subsection{Evaluation on ISPG generation-1 content}

Using Table~\ref{tab:gen1} with the left-Weyl conjugate basis
$\{L_L,\,Q_L,\,e_L^c,\,u_L^c,\,d_L^c\}$
(doublets counted once per $\mathrm{SU}(2)$ state, colored
fields weighted by $N_c$):

\begin{center}
\begin{tabular}{clcc}
\toprule
\# & Anomaly & Value & Status \\
\midrule
1 & $\mathrm{U}(1)_Y^3$           & $0$ & pass\\
2 & $\mathrm{SU}(2)^2\mathrm{U}(1)_Y$  & $0$ & pass\\
3 & $\mathrm{SU}(3)^2\mathrm{U}(1)_Y$  & $0$ & pass\\
4 & grav${}\times\mathrm{U}(1)_Y$ & $0$ & pass\\
5 & Witten $\mathrm{SU}(2)$       & $4$ doublets (even) & pass\\
\bottomrule
\end{tabular}
\end{center}

All five anomalies cancel exactly (verified with rational
arithmetic in the script).

\subsection{What this tells us and what it does not}

\fromSM\ \numerical\ The anomaly conditions are exact
consistency checks on the imported left-Weyl SM hypercharge
assignment. With the SM representation content and the standard
hypercharge normalization fixed, changing an individual entry in
$(Y_{L_L},Y_{Q_L},Y_{e^c},Y_{u^c},Y_{d^c})$ generically
reactivates at least one anomaly. The specific SM values
$(Y_{L_L},Y_{Q_L},Y_{e^c},Y_{u^c},Y_{d^c})
=(-1,+\tfrac13,+2,-\tfrac43,+\tfrac23)$ are imported through
the Gen-1 formula \eqref{eq:gen1formula} combined with
Gell-Mann--Nishijima; they are not independently derived from
archion dynamics in this paper.

The integer pattern~\eqref{eq:YNc_integer} is therefore a
parallel charge-quantization / integer-grading clue, not the
theorem that cancels anomalies. The five anomaly equations above
are what cancel for the imported assignment. In particular,
the anomaly sums are homogeneous in the hypercharges, so a common
rescaling can preserve cancellation while changing the integer
normalization of $Y\,N_c$; conversely, an arbitrary integer
$Y\,N_c$ pattern need not satisfy the five sums. What Phase E
therefore establishes is: the SM hypercharge assignment, once
imported onto the ISPG generation-1 content through the Gen-1
map, passes all five standard anomaly tests. In the Pati--Salam
unification it is known that $Y$ and $N_c$ are linked so that
$Y\,N_c$ is integer; ISPG sees the same integrality as a
compatibility clue, without postulating a larger unifying gauge
group.

\subsection{Phase E verdict}

All five standard anomaly conditions are satisfied by the
imported SM hypercharge assignment on the ISPG generation-1
content. This is a consistency check of the imported charges
against the cavity's integer grading; it does not on its own
constitute a derivation of the SM hypercharges from archion
dynamics.

\section{Phase F --- Generation replication}
\label{sec:phaseF}

\paragraph{Script.} \texttt{python/phase\_f\_generation\_replication.py}.

\subsection{Structural identity of the three generations}

\fromSM\ The three SM generations carry identical
$(Q,T_3,Y,N_c)$ assignments; they differ only by mass and by
entries of the CKM and PMNS mixing matrices. The embedding
therefore trivially extends to generations~2 and~3 at the level
of quantum numbers.

\subsection{Anomaly cancellation for $n_{\rm gen}\geq 2$}

\derived\ Each of the five anomaly conditions is linear in the
generation count: if $A_i^{(1)}=0$ on generation 1, then
$A_i^{(n)}=n\,A_i^{(1)}=0$ for any $n\geq 1$. The script
verifies this explicitly for $n=1,2,3$.

\subsection{The number three: three-axis hypothesis}

\modeldep\ The SM takes the number of generations as an empirical
input (the LEP measurement $N_\nu=2.984\pm 0.008$ from the
$Z$-boson width \cite{PDG2024}). ISPG offers a geometric
reading:
\begin{itemize}
  \item space is three-dimensional;
  \item a 3D archion admits three orthogonal vibration axes;
  \item each axis can host one observed family of cavity modes;
  \item the observed three generations can therefore be mapped
    to the three axes.
\end{itemize}
Under this reading, the correspondence
\begin{equation}
  \text{Gen-1}\leftrightarrow\hat{x},\quad
  \text{Gen-2}\leftrightarrow\hat{y},\quad
  \text{Gen-3}\leftrightarrow\hat{z}
\end{equation}
ties each observed generation to one spatial axis. The exact
result in this phase is the quantum-number replication; the
three-axis origin of $n_{\rm gen}=3$ remains model-dependent
until the cavity dynamics also exclude additional family-like
modes. In the companion MASS work (\S\ref{sec:prereq}), the
charged-lepton Koide geometry and Sargent-type weak-decay
scaling show strong numerical compatibility with this reading,
whereas the individual labels $N_\mu=72$, $N_\tau=295$, and the
ratio $N_\tau/N_e=59$ remain selected labels whose exact
3D-cavity derivation is still open.

\subsection{Fermion count}

\derived\ The total number of left-handed Weyl fermion states
(counting colors) in three generations is
\begin{equation}
  3\times\bigl[\,2\cdot 1 + 2\cdot 3 + 1 + 3 + 3\,\bigr]
  \;=\; 3\times 15 \;=\; 45,
\end{equation}
matching the SM total exactly. Every one of the $45$ states
satisfies~\eqref{eq:GMN}.

\subsection{Phase F verdict}

The three generations exactly replicate generation 1 at the
level of imported SM quantum numbers. If the three-axis
hypothesis is adopted, this replica structure can be assigned
to the three orthogonal spatial axes. Anomaly cancellation is
automatic for any number of replicated generations. Thus Phase
F closes the generation-replication bookkeeping, while the
claim that spatial dimensionality dynamically fixes
$n_{\rm gen}=3$ remains a model-dependent open step.

\section{Synthesis}
\label{sec:synthesis}

\subsection{Cross-sector mass-law compatibility}
\label{sec:unified}

\numerical\ Applying the large-$N$ form
$m_N=(N/N_e)^{2}\,m_e$ with $N_e=5$ uniformly across the twelve
listed non-neutrino massive Standard-Model states (three
charged leptons, six quarks, $W$, $Z$, $H$) yields the
following compatibility table.
Script: \texttt{python/higgs\_replacement\_unified.py}.

\begin{center}
\begin{tabular}{llrrrr}
\toprule
Particle & Sector & $N$ &
 $m_{\text{PDG}}$ (MeV) & $m_{\text{ISPG}}$ (MeV) & residual \\
\midrule
$e$       & lepton &    $5$ &      $0.5110$ &      $0.5110$ & $+0.00\%$ \\
$\mu$     & lepton &   $72$ &    $105.66$   &    $105.96$   & $+0.29\%$ \\
$\tau$    & lepton &  $295$ &   $1776.86$   &   $1778.79$   & $+0.11\%$ \\
$u$       & quark  &   $10$ &      $2.16$   &      $2.044$  & $-5.37\%$ \\
$d$       & quark  &   $15$ &      $4.67$   &      $4.599$  & $-1.52\%$ \\
$s$       & quark  &   $67$ &     $93.40$   &     $91.76$   & $-1.76\%$ \\
$c$       & quark  &  $249$ &   $1270.0$    &   $1267.3$    & $-0.21\%$ \\
$b$       & quark  &  $452$ &   $4180.0$    &   $4176.0$    & $-0.10\%$ \\
$t$       & quark  & $2906$ & $172\,570$    & $172\,612$    & $+0.02\%$ \\
$W$       & gauge  & $1986$ &  $80\,369$    &  $80\,619$    & $+0.31\%$ \\
$Z$       & gauge  & $2115$ &  $91\,188$    &  $91\,433$    & $+0.27\%$ \\
$H$       & Higgs  & $2479$ & $125\,250$    & $125\,613$    & $+0.29\%$ \\
\bottomrule
\end{tabular}
\end{center}

Sector-wise absolute residuals: leptons $\leq 0.29\%$, gauge
bosons $\leq 0.31\%$, Higgs $0.29\%$, quarks $\leq 5.37\%$
(current-mass scheme). Overall mean absolute residual is
$0.85\%$; the single outlier is the up-quark current mass
(scheme-dependent; the pole/$\overline{\rm MS}$ discrepancy
at this scale is itself of order several percent).

As a further consistency check on the same ladder, the ratio
\[
  \left(\frac{N_W}{N_Z}\right)^{\!2}=0.88173
  \quad\text{vs.}\quad
  \left(\frac{m_W}{m_Z}\right)_{\text{PDG}}=0.88135,
\]
i.e.\ $0.04\%$ agreement with the SM on-shell identity
$m_W/m_Z=\cos\theta_W^{\rm os}$; this is
Eq.~\eqref{eq:cosThetaW} re-evaluated on the full table. The
$W$ mass used here is the PDG 2025 average; the newer CMS 2026
value $m_W=80\,360.2\pm 9.9\,\mathrm{MeV}$ does not materially
change the compatibility status.

\paragraph{Status of this table: compatibility, not derivation.}
Every $N_f$ is fit individually to the observed mass through
\eqref{eq:massN}; the table is therefore a demonstration that a
single integer ladder with a single anchor $(N_e=5,\,m_e)$
accommodates the twelve listed non-neutrino massive SM states
within sub-percent residuals across the heavy sector. Neutrino
masses are not included in this ladder table and remain a
separate open sector. The structural content that survives
parameter counting is the \emph{uniformity}: one formula, one
anchor, one ladder, twelve listed particles across four sectors,
no Yukawa coupling, no scalar VEV, no sector-specific coupling
constant. Promotion from fit to derivation requires
producing the $N_f$ themselves from cavity dynamics (items E1
and E5 of \S\ref{sec:extension}); the partial 3D-cavity result
for $N_e=5$ and for
$\kappa_\tau^{2}/\kappa_e^{2}\simeq N_\tau$ recorded under E5
is the first step of that programme.

\subsection{MVP closure}

Combining phases A--F:

\begin{center}
\begin{tabular}{clc}
\toprule
Phase & Content & Status \\
\midrule
A & gauge sector $\mathrm{SU}(3)\times\mathrm{SU}(2)\times\mathrm{U}(1)$
  & SU(3) octet count; E6 algebra/vertices deferred; EW numerical \\
B/C & Gen-1 fermion quantum numbers $(Q,T_3,Y,N_c)$
  & imported-SM compatibility \\
D & EW breaking without Higgs & structural candidate (E7 open) \\
E & five gauge anomalies
  & exact imported-hypercharge check \\
F & three generations as three axes
  & compatibility (model-dep.)\\
\bottomrule
\end{tabular}
\end{center}

The ISPG framework therefore \emph{consistently hosts} the SM
content at the level checked here: none of the gauge-theoretic
requirements of the SM (gauge group dimension, fermion
content, anomaly cancellation, $m_W/m_Z$ consistency,
generation count) is in contradiction with the ISPG archion +
Mathieu-ladder + three-axis framework. A \emph{complete}
structural embedding would further require (i) the
$\mathrm{SU}(3)$ Lie-algebra structure beyond the mode count
(E6), (ii) a derived archion-interior mechanism for gauge-boson
mass generation (E7), (iii) the individual cavity labels
$N_W,N_Z$ (E1), and (iv) a cavity origin for chirality (E2).
The present paper reduces the SM-compatibility question to
this well-defined list of open items; it does not close them.

\subsection{What is derived versus what is imported versus what
is fit}

\begin{itemize}
  \item \derived\ (structurally, from ISPG substrate):
    the Mathieu mass law $m\propto N^2$ (from the companion
    MASS paper);
    the linearity of gauge anomalies in $n_{\rm gen}$.
  \item \numerical\ (partial compatibility clue):
    the \emph{octet count} of confined modes in the
    $D_3$-deformed oscillon toy/robustness analysis; the
    full $\mathrm{SU}(3)_c$ algebra, vertices, and
    $G_2(X)$ parameters remain deferred to E6; the checked
    integer pattern $Y\,N_c\in\mathbb{Z}$ in the imported SM
    assignment.
  \item \numerical/fit\ (diagnostic, not predictive):
    the generation-1 formula~\eqref{eq:gen1formula} (three
    ingredients vs three data points); the individual
    cavity labels $N_W,N_Z,N_H$ (fit to $m_W,m_Z,m_H$); the
    $m_W/m_Z$ ratio test (self-consistent given the
    individual fits).
  \item \fromSM\ (imported without re-derivation):
    the gauge-anomaly formulae (1)--(5); the
    Gell-Mann--Nishijima relation $Q=T_3+Y/2$; the
    hypercharge values $Y_i$ (via Gen-1 formula +
    Gell-Mann--Nishijima); the chirality assignment
    (left-Weyl basis); the identical $(Q,T_3,Y,N_c)$
    across generations.
  \item \modeldep\ (hypothesis):
    the three-axis reading of the three generations;
    the archion-interior mechanism for gauge-boson mass.
\end{itemize}
This separation makes explicit that the ``structural
embedding'' claim rests on items in the first category;
items in the second and third categories must be promoted
to the first before the embedding becomes a derivation.

\subsection{Relation to existing ISPG papers}

This article does not introduce any new field, coupling, or
action term: the action~\eqref{eq:action} is unchanged from
\cite{Lotishvili2026,Lotishvili2026Q}. The Mathieu ladder is
unchanged from the companion MASS paper
(\texttt{MASS/ISPG\_FrequencyToMass.tex}). The
eight-gluon derivation is unchanged from the dedicated
$D_3$-oscillon analysis
(\texttt{SM/ISPG\_EightGluon\_Robustness.md}). The present
paper \emph{synthesizes} these earlier results into a single
embedding statement, with the quantum-number maps of
\S\ref{sec:phaseBC} and the anomaly analysis of
\S\ref{sec:phaseE} as new content.

\section{Extension program}
\label{sec:extension}

The following questions lie beyond the MVP embedding. We list
them not as weaknesses of ISPG but as well-posed structural
extensions, each tied to an identifiable computational or
theoretical task.

\paragraph{E1. Individual cavity labels $N_W, N_Z$.}
The ratio consistency check is already recorded in
\S\ref{sec:phaseA_ew}; it becomes an independent prediction
only if the individual labels $N_W$ and $N_Z$ are derived
from the vector-cavity spectrum.  That derivation requires
diagonalizing the vector Laplacian on the 3D archion
background. A plausible attribution (to be tested numerically)
is $N_W$ to the $\ell=1$ triplet with a
specific radial quantum number and $N_Z$ to a radial excitation
obtained by orthogonalization against a hypercharge-like $\ell=0$
mode. This is a finite-dimensional numerical task.

\paragraph{E2. Chirality from the cavity.}
The left/right distinction is currently imported from the Dirac
structure of $S_m$. A cavity origin would require posing the
chiral eigenvalue problem on the archion background: identify a
parity-breaking boundary condition or a topological invariant
that splits the spectrum into left-handed weak doublets and
right-handed weak singlets with the imported SM quantum numbers.
The success criterion is to recover the SM chirality assignment
without postulating it. This is a conceptually distinct task from
the embedding itself.

\paragraph{E3. CKM and PMNS matrices.}
Under the three-axis hypothesis the three generations are assigned
to a spatial basis, but the quark and lepton mixing sectors are
not the same problem. The small CKM mixing could be represented by
perturbative non-orthogonality of the quark axes, while the large
PMNS mixing requires a separate dynamical mechanism in the
neutrino sector. A first-principles derivation must recover the
CKM Wolfenstein parameters, the PMNS angles, the CKM and PMNS CP
phases, the neutrino mass ordering, and the Majorana/Dirac
character of neutrinos from cavity quantities. At present these
are open targets.

\paragraph{E4. $\alpha_{\rm NL}$ and $\Phi_c$.}
The ISPG Lagrangian contains two dimensionful scales
($\alpha_{\rm NL}$ for the nonlinear self-interaction, $\Phi_c$
for the archion amplitude) that
are currently fit. A first-principles derivation would proceed
via the nonlinear archion profile at fixed cosmological
background; recent work identifies this as a BVP in the
stationary $\Phi$ equation. The notation $\alpha_{\rm NL}$ is used
to keep this nonlinear ISPG scale distinct from the fine-structure
constant $\alpha_{\rm EM}$.

\paragraph{E5. Exact mass hierarchy.}
The charged-lepton hierarchy ($m_e,m_\mu,m_\tau$) is numerically
reproduced by the three-axis hypothesis to the precisions cited
in \S\ref{sec:prereq}; the quark hierarchy
($m_u\ll m_c\ll m_t$) is at the same heuristic level. Promoting
the three-axis ansatz to a deduction from the archion ground
state is the central open task for a mass-exact (as opposed to
compatibility-level) embedding.

\numerical\ A partial structural result is available for the
anchor label itself: the electron index $N_e=5$ matches the
$\ell=2$ degeneracy count of the 3D spherical cavity (five
independent spherical harmonics $Y_2^{m}$, $m=-2,\ldots,+2$).
In the same 3D cavity analysis the ratio $N_\tau^{\,2}/N_e^{\,2}$
is not reproduced. The actual partial result is
$\kappa_\tau^{2}/\kappa_e^{2}\simeq 292.43$, which matches the
selected label $N_\tau=295$ to within $0.9\%$; equivalently,
$\kappa_\tau^{2}/(5\kappa_e^{2})\simeq 58.49$ is close to the
selected ratio $N_\tau/N_e=59$. These results fix part of the
kinematic skeleton of the three-axis ansatz (dimensionality of
the mode space, radial wavenumber ratios) at the
cavity-geometry level, but they do not yet derive the exact
integer labels.

\deferred\ Extending the partial result to absorb the $0.9\%$
residual requires a Floquet treatment of the periodic archion:
the present analysis uses a peak-amplitude approximation of the
nonlinear potential, whereas the correct procedure replaces this
by the time-averaged effective potential computed from the full
Floquet spectrum at deformation $q\simeq 1.853$. This is a
well-defined computational task whose completion would promote
E5 from compatibility into derivation.

\paragraph{E6. $\mathrm{SU}(3)$ Lie-algebra structure beyond
the count.} Extracting the Gell-Mann commutation relations
$[T^a,T^b]=if^{abc}T^c$ with the standard $\mathrm{SU}(3)$
structure constants from the trilinear mode-coupling vertices
of the eight confined modes of the $D_3$-deformed oscillon.
Without this step, the gluon count is consistent with several
Lie algebras of dimension $8$, not only $\mathrm{SU}(3)$.
\numerical\ A partial group-theoretic result is available: if
the three spatial axes are identified with the fundamental
$\mathbf{3}$ of color (as in the three-axis hypothesis used for
the lepton sector), the eight generators of unitary mixings
among them close on the $\mathrm{SU}(3)$ Lie algebra with the
textbook structure constants $f^{abc}$ recovered to machine
precision ($\leq 2\times 10^{-16}$ on all nine independent
entries, full antisymmetry, and the Jacobi identity;
script \texttt{python/gluon\_su3\_closure.py}). This fixes
the target algebra but does not replace the dynamical
extraction of $f^{abc}$ from the oscillon vertices, which is
the content of E6 proper.

\paragraph{E7. Archion-interior gauge-boson mass generation.}
A derivation showing that an $\mathrm{SU}(2)\times\mathrm{U}(1)$
gauge field minimally coupled to the archion acquires a mass
matrix matching $(M_W,M_Z)$ while leaving the photon massless.
This is the concrete replacement for the SM Higgs mechanism
that the ISPG action forbids; its absence is the single largest
dynamical gap in the present paper. The operational target is:
write the imported $\mathrm{SU}(2)_L\times\mathrm{U}(1)_Y$ gauge
action on the archion background, specify its coupling to
$\Phi(r)$, derive the mass matrix in the $(W^1,W^2,W^3,B)$ basis,
show that one eigenvalue remains exactly zero (the photon mode),
and check the gauge-consistency conditions needed for current
conservation and unitarity. Only then can the observed $M_W$,
$M_Z$, and the E1 cavity labels be treated as derived rather than
fit. Plausible route: expand the gauge kinetic term around the
radially varying archion profile $\Phi(r)$ and show that the
volume integrals yield the observed mass pattern when combined
with the cavity labels from E1.

\medskip
None of these items conflicts with the compatibility
established in \S\S\ref{sec:phaseA}--\ref{sec:phaseF}. Each is
a well-defined computational or theoretical task whose
resolution would promote the present paper from
\emph{compatibility check} to \emph{derivation}. Until E1, E6,
and E7 are closed, the electroweak and strong sectors should
not be represented as ``derived'' from ISPG but as
``consistently hosted in'' ISPG.

\section{Conclusion}
\label{sec:conclusion}

The SM gauge and fermion content is compatible with the ISPG
cavity framework at the level checked here: the eight-gluon
count matches $\dim\mathrm{SU}(3)_{\mathrm{adj}}$, the Mathieu
ladder hosts the electroweak bosons consistently, the SM
hypercharges pass all five anomaly tests when imported via
the generation-1 fit formula combined with
Gell-Mann--Nishijima, and three generations admit a
model-dependent geometric reading in terms of three spatial
axes. Neutrinos are explicitly excluded from the Mathieu-ladder
fermion formula and are handled in the companion MASS paper as
polarization-rotation modes.

The embedding is however \emph{partial}: the $\mathrm{SU}(3)$
Lie algebra beyond the count (E6), the archion-interior mass
mechanism replacing the SM Higgs (E7), the individual cavity
labels $N_W,N_Z$ (E1), and the cavity origin of chirality (E2)
are all outstanding. The $\cos\theta_W^{\rm os}$ test of
\S\ref{sec:phaseA_ew} is a self-consistency check rather than
a derivation.

The value of the present paper lies in having separated these
items from one another and from the ISPG substrate: each open
item is now a well-posed computational or theoretical task,
and none requires modifying the ISPG action~\eqref{eq:action}.
Closing them in sequence would convert the present
compatibility result into a full derivation of the SM from the
single-scalar substrate; the structural path is laid out in
\S\ref{sec:extension}.

\appendix

\section{Python verification scripts}
\label{app:scripts}

All scripts are self-contained and run from
\texttt{python/} with no external state.

\begin{center}
\begin{tabular}{lp{0.55\textwidth}l}
\toprule
Script & Purpose & Phase \\
\midrule
\texttt{phase\_a\_gauge\_sector.py}
  & EW boson Mathieu indices; on-shell $\cos\theta_W$ ratio test
  & A \\
\texttt{phase\_c\_quantum\_number\_map.py}
  & Gen-1 formula $N=5[1+N_c(1-|Q|)]$ on $(e,\nu_e,u,d)$
  & C.1 \\
\texttt{phase\_c\_step2\_T3\_Y\_map.py}
  & Gell-Mann--Nishijima for all 7 L+R Gen-1 fermions; $Y\,N_c$ integrality
  & C.2 \\
\texttt{phase\_d\_ew\_breaking.py}
  & $m\propto N^2$ test on $W,Z,H$; absence of $\phi^4$ in the action
  & D \\
\texttt{phase\_e\_anomalies.py}
  & Five anomaly conditions, exact rational arithmetic
  & E \\
\texttt{phase\_f\_generation\_replication.py}
  & Structural identity of Gen 1,2,3; anomaly linearity in $n_{\rm gen}$
  & F \\
\texttt{gluon\_su3\_closure.py}
  & Target-algebra compatibility on assumed three colour axes;
    textbook $f^{abc}$ matched to $\leq 2\times 10^{-16}$, but no
    dynamic extraction from confined modes
  & A compatibility; E6 target only \\
\texttt{higgs\_replacement\_unified.py}
  & Unified compatibility table $m=(N/N_e)^2 m_e$ across the 12 listed
    non-neutrino massive SM states (charged leptons, quarks, $W,Z,H$);
    mean residual $0.85\%$
  & A--D \\
\bottomrule
\end{tabular}
\end{center}

Each script prints a human-readable report ending in either a
pass line or an explicit failure indicator. The cumulative
output constitutes a machine-verifiable audit of the embedding.

\clearpage
\phantomsection
\section*{References}
% Consolidated bibliography for standalone and master builds
\begin{thebibliography}{99}

\bibitem{Einstein1915}
A.~Einstein,
``Die Feldgleichungen der Gravitation,''
Sitzungsber.\ Preuss.\ Akad.\ Wiss.\ Berlin, 844 (1915).

\bibitem{Will2014}
C.~M.~Will,
``The Confrontation between General Relativity and Experiment,''
Living Rev.\ Relativ.\ \textbf{17}, 4 (2014).

\bibitem{Will2018}
C.~M.~Will,
\textit{Theory and Experiment in Gravitational Physics},
2nd ed.\ (Cambridge University Press, 2018).

\bibitem{LIGO2016}
B.~P.~Abbott \textit{et al.}\ (LIGO/Virgo),
``Observation of Gravitational Waves from a Binary Black Hole Merger,''
Phys.\ Rev.\ Lett.\ \textbf{116}, 061102 (2016).

\bibitem{LIGO_GW170817}
B.~P.~Abbott \textit{et al.}\ (LIGO/Virgo/Fermi/INTEGRAL),
``Gravitational Waves and Gamma-Rays from a Binary Neutron Star Merger: GW170817 and GRB~170817A,''
Astrophys.\ J.\ Lett.\ \textbf{848}, L13 (2017).

\bibitem{TullyFisher1977}
R.~B.~Tully and J.~R.~Fisher,
``A New Method of Determining Distances to Galaxies,''
Astron.\ Astrophys.\ \textbf{54}, 661 (1977).

\bibitem{McGaugh2016}
S.~S.~McGaugh, F.~Lelli, and J.~M.~Schombert,
``Radial Acceleration Relation in Rotationally Supported Galaxies,''
Phys.\ Rev.\ Lett.\ \textbf{117}, 201101 (2016).

\bibitem{Milgrom1983a}
M.~Milgrom,
``A Modification of the Newtonian Dynamics as a Possible Alternative to the Hidden Mass Hypothesis,''
Astrophys.\ J.\ \textbf{270}, 365 (1983).

\bibitem{Milgrom1999}
M.~Milgrom,
``The Modified Dynamics as a Vacuum Effect,''
Phys.\ Lett.\ A \textbf{253}, 273 (1999).

\bibitem{Famaey2012}
B.~Famaey and S.~S.~McGaugh,
``Modified Newtonian Dynamics (MOND): Observational Phenomenology and Relativistic Extensions,''
Living Rev.\ Relativ.\ \textbf{15}, 10 (2012).

\bibitem{Angus2008}
G.~W.~Angus, B.~Famaey, and H.~S.~Zhao,
``Can MOND take a bullet? Analytical comparisons of three versions of MOND beyond spherical symmetry,''
Mon.\ Not.\ R.\ Astron.\ Soc.\ \textbf{371}, 138 (2006).

\bibitem{Bekenstein2004}
J.~D.~Bekenstein,
``Relativistic Gravitation Theory for the Modified Newtonian Dynamics Paradigm,''
Phys.\ Rev.\ D \textbf{70}, 083509 (2004).

\bibitem{Sotiriou2010}
T.~P.~Sotiriou and V.~Faraoni,
``$f(R)$ Theories of Gravity,''
Rev.\ Mod.\ Phys.\ \textbf{82}, 451 (2010).

\bibitem{BrunetonEsposito2007}
J.-P.~Bruneton and G.~Esposito-Far\`ese,
``Field-Theoretical Formulations of MOND-like Gravity,''
Phys.\ Rev.\ D \textbf{76}, 124012 (2007).

\bibitem{Penrose1965}
R.~Penrose,
``Gravitational Collapse and Space-Time Singularities,''
Phys.\ Rev.\ Lett.\ \textbf{14}, 57 (1965).

\bibitem{Hawking1970}
S.~W.~Hawking and R.~Penrose,
``The Singularities of Gravitational Collapse and Cosmology,''
Proc.\ R.\ Soc.\ Lond.\ A \textbf{314}, 529 (1970).

\bibitem{DeWitt1967}
B.~S.~DeWitt,
``Quantum Theory of Gravity. I. The Canonical Theory,''
Phys.\ Rev.\ \textbf{160}, 1113 (1967).

\bibitem{tHooft1974}
G.~'t~Hooft and M.~Veltman,
``One-Loop Divergences in the Theory of Gravitation,''
Ann.\ Inst.\ Henri Poincar\'e A \textbf{20}, 69 (1974).

\bibitem{BransDicke1961}
C.~Brans and R.~H.~Dicke,
``Mach's Principle and a Relativistic Theory of Gravitation,''
Phys.\ Rev.\ \textbf{124}, 925 (1961).

\bibitem{Horndeski1974}
G.~W.~Horndeski,
``Second-Order Scalar-Tensor Field Equations in a Four-Dimensional Space,''
Int.\ J.\ Theor.\ Phys.\ \textbf{10}, 363 (1974).

\bibitem{BelliniSawicki2014}
E.~Bellini and I.~Sawicki,
``Maximal Freedom at Minimum Cost: Linear Large-Scale Structure in General Modifications of Gravity,''
J.\ Cosmol.\ Astropart.\ Phys.\ \textbf{2014}(07), 050.

\bibitem{Bertotti2003}
B.~Bertotti, L.~Iess, and P.~Tortora,
``A Test of General Relativity Using Radio Links with the Cassini Spacecraft,''
Nature \textbf{425}, 374 (2003).

\bibitem{Williams2004}
J.~G.~Williams, S.~G.~Turyshev, and D.~H.~Boggs,
``Progress in Lunar Laser Ranging Tests of Relativistic Gravity,''
Phys.\ Rev.\ Lett.\ \textbf{93}, 261101 (2004).

\bibitem{Everitt2011}
C.~W.~F.~Everitt \textit{et al.},
``Gravity Probe B: Final Results of a Space Experiment to Test General Relativity,''
Phys.\ Rev.\ Lett.\ \textbf{106}, 221101 (2011).

\bibitem{Shapiro1964}
I.~I.~Shapiro,
``Fourth Test of General Relativity,''
Phys.\ Rev.\ Lett.\ \textbf{13}, 789 (1964).

\bibitem{Blanchet1993}
L.~Blanchet and G.~Sch\"afer,
``Gravitational Wave Tails and Binary Star Systems,''
Classical Quantum Gravity \textbf{10}, 2699 (1993).

\bibitem{BodennerWill2003}
J.~Bodenner and C.~M.~Will,
``Deflection of light to second order: A tool for illustrating principles of general relativity,''
Am.\ J.\ Phys.\ \textbf{71}, 770--773 (2003).

\bibitem{EpsteinShapiro1980}
R.~Epstein and I.~I.~Shapiro,
``Post-post-Newtonian deflection of light by the Sun,''
Phys.\ Rev.\ D \textbf{22}, 2947--2949 (1980).

\bibitem{KeetonPetters2005}
C.~R.~Keeton and A.~O.~Petters,
``Formalism for testing theories of gravity using lensing by compact objects,''
Phys.\ Rev.\ D \textbf{72}, 104006 (2005).

\bibitem{SKA2015}
M.~Kramer \textit{et al.},
``Strong-Field Tests of Gravity Using Pulsars and Black Holes,''
in \textit{Advancing Astrophysics with the Square Kilometre Array} (2015).

\bibitem{Carroll2003}
S.~M.~Carroll, M.~Hoffman, and M.~Trodden,
``Can the Dark Energy Equation-of-State Parameter $w$ Be Less Than $-1$?''
Phys.\ Rev.\ D \textbf{68}, 023509 (2003).

\bibitem{Cline2004}
J.~M.~Cline, S.~Jeon, and G.~D.~Moore,
``The Phantom Menaced: Constraints on Low-Energy Effective Ghosts,''
Phys.\ Rev.\ D \textbf{70}, 043543 (2004).

\bibitem{Sbisa2015}
F.~Sbis\`a,
``Classical and Quantum Ghosts,''
Eur.\ J.\ Phys.\ \textbf{36}, 015009 (2015).

\bibitem{Papallo2018}
G.~Papallo and H.~S.~Reall,
``On the Local Well-Posedness of Lovelock and Horndeski Theories,''
Phys.\ Rev.\ D \textbf{96}, 044019 (2017).

\bibitem{Kovacs2020}
A.~D.~Kov\'acs and H.~S.~Reall,
``Well-Posed Formulation of Scalar-Tensor Effective Field Theory,''
Phys.\ Rev.\ Lett.\ \textbf{124}, 221101 (2020).

\bibitem{Afshordi2007}
N.~Afshordi, D.~J.~H.~Chung, and G.~Geshnizjani,
``Cuscuton: A Causal Field Theory with an Infinite Speed of Sound,''
Phys.\ Rev.\ D \textbf{75}, 083513 (2007).

\bibitem{Chamseddine2013}
A.~H.~Chamseddine and V.~Mukhanov,
``Mimetic Dark Matter,''
J.\ High Energy Phys.\ \textbf{2013}(11), 135.

\bibitem{ArkaniHamed2004}
N.~Arkani-Hamed, H.-C.~Cheng, M.~A.~Luty, and S.~Mukohyama,
``Ghost Condensation and a Consistent Infrared Modification of Gravity,''
J.\ High Energy Phys.\ \textbf{2004}(05), 074.

\bibitem{SageManifolds}
E.~Gourgoulhon, M.~Bejger, and others,
``SageManifolds: Differential Geometry and Tensor Calculus with SageMath,''
\url{https://sagemanifolds.obspm.fr/} (2023).

\bibitem{EHT2022}
Event Horizon Telescope Collaboration,
``First Sagittarius~A* Event Horizon Telescope Results. I.,''
Astrophys.\ J.\ Lett.\ \textbf{930}, L12 (2022).

\bibitem{EHT2019M87}
Event Horizon Telescope Collaboration,
``First M87 Event Horizon Telescope Results. I. The Shadow of the Supermassive Black Hole,''
Astrophys.\ J.\ Lett.\ \textbf{875}, L1 (2019),
DOI: 10.3847/2041-8213/ab0ec7.

\bibitem{EHT2024M87}
Event Horizon Telescope Collaboration,
``The persistent shadow of the supermassive black hole of M87. I. Observations, calibration, imaging, and analysis,''
Astron.\ Astrophys.\ \textbf{681}, A79 (2024),
DOI: 10.1051/0004-6361/202347932.

\bibitem{Cardoso2019}
V.~Cardoso and P.~Pani,
``Testing the Nature of Dark Compact Objects: A Status Report,''
Living Rev.\ Relativ.\ \textbf{22}, 4 (2019).

\bibitem{Miani2023}
A.~Miani \emph{et al.},
``Constraints on the amplitude of gravitational wave echoes from black hole ringdown using minimal assumptions,''
Phys.\ Rev.\ D \textbf{108}, 064018 (2023),
DOI: 10.1103/PhysRevD.108.064018.

\bibitem{LVK2025TGR}
R.~Abbott \emph{et al.}\ (LIGO Scientific, Virgo, and KAGRA Collaborations),
``Tests of general relativity with GWTC-3,''
Phys.\ Rev.\ D \textbf{112}, 084080 (2025),
DOI: 10.1103/PhysRevD.112.084080.

\bibitem{Gogberashvili2018}
M.~Gogberashvili and D.~Singleton,
``Quantum particles trapped in a Schwarzschild black hole can escape in finite time,''
Phys.\ Lett.\ B \textbf{784}, 17--21 (2018).

\bibitem{Gogberashvili2019}
M.~Gogberashvili, A.~Sakharov, and E.~Sarkisyan-Grinbaum,
``Black holes as inner boundaries of the universe,''
Int.\ J.\ Mod.\ Phys.\ D \textbf{28}, 1950178 (2019).

\bibitem{McGaugh2011}
S.~S.~McGaugh,
``Novel Test of Modified Newtonian Dynamics with Gas Rich Galaxies,''
Phys.\ Rev.\ Lett.\ \textbf{106}, 121303 (2011).

\bibitem{Lotishvili2026}
G.~Lotishvili,
``Bi-Conformal Scalar--Tensor Gravity: Exact PPN Equivalence with General Relativity and Falsifiable Predictions,''
Zenodo (2026).

\bibitem{Lotishvili2026MOND}
G.~Lotishvili,
``Deriving MOND from Bi-Conformal Gravity: Vortex Closure, the Interpolating Function, and the Dark-Energy Connection,''
Zenodo (2026).

\bibitem{Lotishvili2026Q}
G.~Lotishvili,
``Quantum Extension of Bi-Conformal Scalar--Tensor Gravity: Canonical Quantization, Particle Spectrum, and Falsifiable Predictions,''
Zenodo (2026).

\bibitem{Hawking1976}
S.~W.~Hawking,
``Breakdown of Predictability in Gravitational Collapse,''
Phys.\ Rev.\ D \textbf{14}, 2460 (1976).

\bibitem{Kibble1976}
T.~W.~B.~Kibble,
``Topology of Cosmic Domains and Strings,''
J.\ Phys.\ A \textbf{9}, 1387 (1976).

\bibitem{DESI2024}
DESI Collaboration,
``DESI 2024 VI: Cosmological constraints from the measurements
of baryon acoustic oscillations,''
Phys.\ Rev.\ D \textbf{110} (2024) 023507.

\bibitem{DESIDR2}
DESI Collaboration, A.~Abdul-Karim \textit{et al.},
``DESI DR2 Results II: Measurements of Baryon Acoustic
Oscillations and Cosmological Constraints,''
Phys.\ Rev.\ D \textbf{112} (2025) 083515.

\bibitem{JWST2023}
S.~L.~Finkelstein \textit{et al.},
``A Long Time Ago in a Galaxy Far, Far Away: A Candidate $z \sim 12$ Galaxy,''
Astrophys.\ J.\ Lett.\ \textbf{940}, L55 (2022).

\bibitem{Euclid2024}
Euclid Collaboration,
``Euclid. I. Overview of the Euclid Mission,''
Astron.\ Astrophys.\ \textbf{662}, A112 (2022).

\bibitem{MICROSCOPE2022}
P.~Touboul \textit{et al.}\ (MICROSCOPE Collaboration),
``MICROSCOPE Mission: Final Results of the Test of the Equivalence Principle,''
Phys.\ Rev.\ Lett.\ \textbf{129}, 121102 (2022).

\bibitem{PoundRebka1960}
R.~V.~Pound and G.~A.~Rebka, Jr.,
``Apparent Weight of Photons,''
Phys.\ Rev.\ Lett.\ \textbf{4}, 337 (1960).

\bibitem{VessotLevine1980}
R.~F.~C.~Vessot and M.~W.~Levine,
``A Test of the Equivalence Principle Using a Space-Borne Clock,''
Gen.\ Relativ.\ Gravit.\ \textbf{10}, 181 (1979);
R.~F.~C.~Vessot \textit{et al.},
``Test of Relativistic Gravitation with a Space-Borne Hydrogen Maser,''
Phys.\ Rev.\ Lett.\ \textbf{45}, 2081 (1980).

\bibitem{MichelsonMorley1887}
A.~A.~Michelson and E.~W.~Morley,
``On the Relative Motion of the Earth and the Luminiferous Ether,''
Am.\ J.\ Sci.\ \textbf{34}, 333 (1887).

\bibitem{LambHydro}
H.~Lamb,
\textit{Hydrodynamics},
6th ed.\ (Cambridge University Press, 1932).

\bibitem{Fisher1948}
I.~Z.~Fisher,
``Scalar Mesostatic Field with Regard for Gravitational Effects,''
Zh.\ Eksp.\ Teor.\ Fiz.\ \textbf{18}, 636 (1948)
[arXiv:gr-qc/9911008 (English translation)].

\bibitem{JNW1968}
A.~I.~Janis, E.~T.~Newman, and J.~Winicour,
``Reality of the Schwarzschild Singularity,''
Phys.\ Rev.\ Lett.\ \textbf{20}, 878 (1968).

\bibitem{Bronnikov1973}
K.~A.~Bronnikov,
``Scalar-Tensor Theory and Scalar Charge,''
Acta Phys.\ Pol.\ B \textbf{4}, 251 (1973).

\bibitem{Ellis1973}
H.~G.~Ellis,
``Ether Flow through a Drainhole: A Particle Model in General Relativity,''
J.\ Math.\ Phys.\ \textbf{14}, 104 (1973).

\bibitem{Papapetrou1954}
A.~Papapetrou,
``Eine Theorie des Gravitationsfeldes mit einer Feldfunktion,''
Z.\ Phys.\ \textbf{139}, 518 (1954).

\bibitem{Yilmaz1958}
H.~Yilmaz,
``New Approach to General Relativity,''
Phys.\ Rev.\ \textbf{111}, 1417 (1958).

\bibitem{Boonserm2020}
P.~Boonserm, T.~Ngampitipan, A.~Simpson, and M.~Visser,
``Triple Path to the Exponential Metric,''
Found.\ Phys.\ \textbf{50}, 1368 (2020).

\bibitem{Bronnikov2012}
K.~A.~Bronnikov and S.~G.~Rubin,
``Instabilities of Wormholes and Regular Black Holes Supported by a Phantom Scalar Field,''
Phys.\ Rev.\ D \textbf{86}, 024028 (2012).

\bibitem{LISA2017}
P.~Amaro-Seoane \textit{et al.},
``Laser Interferometer Space Antenna,''
arXiv:1702.00786 [astro-ph.IM] (2017).

\bibitem{Adelberger2003}
E.~G.~Adelberger, B.~R.~Heckel, and A.~E.~Nelson,
``Tests of the Gravitational Inverse-Square Law,''
Annu.\ Rev.\ Nucl.\ Part.\ Sci.\ \textbf{53}, 77 (2003).

\bibitem{Clowe2006}
D.~Clowe, M.~Brada\v{c}, A.~H.~Gonzalez, M.~Markevitch,
S.~W.~Randall, C.~Jones, and D.~Zaritsky,
``A Direct Empirical Proof of the Existence of Dark Matter,''
Astrophys.\ J.\ Lett.\ \textbf{648}, L109 (2006).

\bibitem{Freire2012}
P.~C.~C.~Freire \textit{et al.},
``The relativistic pulsar--white dwarf binary PSR~J1738+0333~II.\ The most stringent test of scalar-tensor gravity,''
Mon.\ Not.\ R.\ Astron.\ Soc.\ \textbf{423}, 3328 (2012).

\bibitem{Archibald2018}
A.~M.~Archibald \textit{et al.},
``Universality of free fall from the orbital motion of a pulsar in a stellar triple system,''
Nature \textbf{559}, 73 (2018).

\bibitem{Voisin2025}
G.~Voisin \textit{et al.},
``Explanation of the exceptionally strong timing noise of PSR~J0337+1715 by a circum-ternary planet and consequences for gravity tests,''
Astron.\ Astrophys.\ \textbf{693}, A143 (2025),
DOI: 10.1051/0004-6361/202452100.

\bibitem{Clarke1993}
C.~J.~S.~Clarke,
\emph{The Analysis of Space-Time Singularities}
(Cambridge University Press, 1993).

\bibitem{Racz2023}
I.~R\'acz,
``Spacetime singularities and curvature blow-ups,''
Gen.\ Relativ.\ Gravit.\ \textbf{55}, 3 (2023),
DOI: 10.1007/s10714-022-03053-9.

\bibitem{Abbott2017GW}
B.~P.~Abbott \emph{et al.} (LIGO Scientific and Virgo Collaborations),
\emph{GW170817: Observation of Gravitational Waves from a Binary Neutron Star Inspiral},
Phys.\ Rev.\ Lett.\ \textbf{119}, 161101 (2017).

\bibitem{SkordisZlosnik2021}
C.~Skordis and T.~Z\l{}o\'snik,
\emph{New relativistic theory for Modified Newtonian Dynamics},
Phys.\ Rev.\ Lett.\ \textbf{127}, 161302 (2021).

\bibitem{Barbour1999}
J.~Barbour,
\emph{The End of Time: The Next Revolution in Physics}
(Oxford University Press, 1999).

\bibitem{ConnesRovelli1994}
A.~Connes and C.~Rovelli,
``Von~Neumann algebra automorphisms and time-thermodynamics relation
in generally covariant quantum theories,''
Class.\ Quantum Grav.\ \textbf{11}, 2899 (1994).

\bibitem{Whitehead1929}
A.~N.~Whitehead,
\emph{Process and Reality}
(Macmillan, 1929).

\bibitem{Penrose2010}
R.~Penrose,
\emph{Cycles of Time: An Extraordinary New View of the Universe}
(Bodley Head, 2010).

% --- Particle / Mathieu / holography supplement (FrequencyToMass, quantum sector) ---

\bibitem{Koide1982}
Y.~Koide,
``Fermion--boson two-body bound systems and the fermion mass hierarchy,''
Lett.\ Nuovo Cim.\ \textbf{34}, 201 (1982).

\bibitem{DLMFMathieu}
NIST Digital Library of Mathematical Functions,
Chap.~28, \textit{Mathieu Functions and Hill's Equation},
\url{https://dlmf.nist.gov/28},
NIST, Gaithersburg, MD (release 1.2.1, 2024).

\bibitem{McLachlan1947}
N.~W.~McLachlan,
\emph{Theory and Application of Mathieu Functions}
(Oxford University Press, 1947).

\bibitem{PDG2024}
R.~L.~Workman \textit{et al.}\ (Particle Data Group),
``Review of Particle Physics,''
Prog.\ Theor.\ Exp.\ Phys.\ \textbf{2024}, 083C01 (2024).

\bibitem{RyuTakayanagi2006}
S.~Ryu and T.~Takayanagi,
``Holographic Derivation of the Entanglement Entropy from AdS/CFT,''
Phys.\ Rev.\ Lett.\ \textbf{96}, 181602 (2006).

\bibitem{NishiokaTakayanagi2009}
T.~Nishioka, S.~Ryu, and T.~Takayanagi,
``Holographic Entanglement Entropy: An Overview,''
J.\ Phys.\ A \textbf{42}, 504008 (2009).

\bibitem{Bombelli1986}
L.~Bombelli, R.~K.~Koul, J.~Lee, and R.~D.~Sorkin,
``A Quantum Source of Entropy for Black Holes,''
Phys.\ Rev.\ D \textbf{34}, 373 (1986).

\bibitem{MaldacenaSusskind2013}
J.~Maldacena and L.~Susskind,
``Cool horizons for entangled black holes,''
Fortschr.\ Phys.\ \textbf{61}, 781 (2013).

\bibitem{Maldacena1998}
J.~M.~Maldacena,
``The Large $N$ limit of superconformal field theories and supergravity,''
Adv.\ Theor.\ Math.\ Phys.\ \textbf{2}, 231 (1998).

\bibitem{Leggett2001}
A.~J.~Leggett,
``Bose--Einstein condensation in the alkali gases: Some fundamental concepts,''
Rev.\ Mod.\ Phys.\ \textbf{73}, 307 (2001).

\bibitem{PethickSmith2008}
C.~J.~Pethick and H.~Smith,
\emph{Bose--Einstein Condensation in Dilute Gases},
2nd ed.\ (Cambridge University Press, 2008).

\bibitem{GleiserWatkins1993}
I.~Gleiser and R.~Watkins,
``Towards a theory of oscillons,''
Phys.\ Rev.\ D \textbf{47}, 5197 (1993).

\bibitem{Copeland1995}
E.~J.~Copeland, M.~Gleiser, and H.-R.~Muller,
``Oscillons: Resonant configurations during damping in one dimension,''
Phys.\ Rev.\ D \textbf{52}, 1920 (1995).

\bibitem{Nicolis2009}
A.~Nicolis, R.~Rattazzi, and E.~Trincherini,
``The Galileon as a local modification of gravity,''
Phys.\ Rev.\ D \textbf{79}, 064036 (2009).

\bibitem{Deffayet2009}
C.~Deffayet, G.~Esposito-Far\`ese, and A.~Vikman,
``Covariantizing Galileons,''
Phys.\ Rev.\ D \textbf{79}, 084003 (2009).

\bibitem{Goldstone1961}
J.~Goldstone,
``Field theories with ``superconductor'' solutions,''
Nuovo Cimento \textbf{19}, 154 (1961).

\bibitem{ColemanWeinberg1973}
S.~Coleman and E.~Weinberg,
``Radiative Corrections as the Origin of Spontaneous Symmetry Breaking,''
Phys.\ Rev.\ D \textbf{7}, 1888 (1973).

\bibitem{Higgs1964}
P.~W.~Higgs,
``Broken Symmetries and the Masses of Gauge Bosons,''
Phys.\ Rev.\ Lett.\ \textbf{13}, 508 (1964).

\bibitem{EnglertBrout1964}
F.~Englert and R.~Brout,
``Broken Symmetry and the Mass of Gauge Vector Mesons,''
Phys.\ Rev.\ Lett.\ \textbf{13}, 321 (1964).

\bibitem{Guralnik1964}
G.~S.~Guralnik, C.~R.~Hagen, and T.~W.~B.~Kibble,
``Global Conservation Laws and Massless Particles,''
Phys.\ Rev.\ Lett.\ \textbf{13}, 585 (1964).

\bibitem{Hopkins2010}
P.~F.~Hopkins \textit{et al.},
``Mergers, Minor Mergers, and the Origin of Disklike Galaxies,''
Astrophys.\ J.\ \textbf{715}, 202 (2010).

\bibitem{Stewart2008}
K.~R.~Stewart \textit{et al.},
``Merger Histories of Galaxy Halos and Implications for Disk Survival,''
Astrophys.\ J.\ \textbf{683}, 597 (2008).

\end{thebibliography}


\end{document}
